% GENERATED FILE. Edit canonical lessons, then run npm run generate:books.
% canonical-source-hash: fnv1a64:15f468e16acbf093
% canonical-lessons: ML-C412-curukkattil, ML-C412-mariccu, ML-C412-athava, ML-C412-potuve, ML-C412-ennuveccal, ML-R412-review-405-412

\chapter{In Short}
\label{ch:ml-in-short}
\hlchaptermodality{412}

\begin{chapteropening}
\textbf{By the end of this chapter:} \emph{I can say in short, on the contrary, or, generally and that means.}

Complete the last lesson of chapter 412: I can say in short, on the contrary, or, generally and that means.
\end{chapteropening}

\section[\texorpdfstring{\ml{ചുരുക്കത്തിൽ} (curukkattil)}{ചുരുക്കത്തിൽ (curukkattil)}]{\ml{ചുരുക്കത്തിൽ} (curukkattil) — in short}

\label{lesson:ML-C412-curukkattil}



\begin{quote}
Before the new one: say the Malayalam for actually, then the Malayalam for in reality.
\end{quote}

\subsection*{You'll want to know: \ml{ചുരുക്കത്തിൽ}}
\textbf{\ml{ചുരുക്കത്തിൽ}} — \emph{curukkattil} — \textquotedblleft{}in short\textquotedblright{}.

\begin{quote}
\ml{ചുരുക്കത്തിൽ}, \ml{യാത്ര} \ml{വളരെ} \ml{നന്നായിരുന്നു}.
\end{quote}

\begin{quote}
\emph{curukkattil, yātra vaḷare nannāyirunnu.}
\end{quote}

\begin{quote}
In short, the trip went very well.
\end{quote}

The locative of \ml{ചുരുക്കം}, shortness.

\begin{cousinweb}
\textbf{\ml{ചുരുക്കത്തിൽ}} (\emph{curukkattil}) is native Dravidian: \ml{ചുരുക്കം}, \textquotedblleft{}contraction, shortness\textquotedblright{} (DEDR 2687), from \ml{ചുരുങ്ങുക}, \textquotedblleft{}to shrink\textquotedblright{}, with the locative ending -\ml{ത്തിൽ}.
\end{cousinweb}

\begin{grammarlens}[title={What you've built}]
One more word for giving reasons and opinions.
\end{grammarlens}

\subsection*{Your turn}
\begin{itemize}
\raggedright
  \item \emph{Say it:} \emph{curukkattil}
  \item \emph{Say it:} \emph{curukkattil}, once more
  \item \emph{Say it:} \emph{curukkattil, yātra vaḷare nannāyirunnu.} — In short, the trip went very well.
  \item \emph{Recall:} say the Malayalam for actually, then the Malayalam for in reality, then say \emph{curukkattil} again
\end{itemize}

\begin{tcolorbox}[breakable,colback=teal!4,colframe=teal!35!black,title={Before you move on}]
What does \ml{ചുരുക്കത്തിൽ} mean? (\textquotedblleft{}In short\textquotedblright{}.) Say it once more.
\end{tcolorbox}
\section[\texorpdfstring{\ml{മറിച്ച്} (maṟiccŭ)}{മറിച്ച് (maṟiccŭ)}]{\ml{മറിച്ച്} (maṟiccŭ) — on the contrary}

\label{lesson:ML-C412-mariccu}



\begin{quote}
Before the new one: say the Malayalam for in reality, then the Malayalam for in short.
\end{quote}

\subsection*{You'll want to know: \ml{മറിച്ച്}}
\textbf{\ml{മറിച്ച്}} — \emph{maṟiccŭ} — \textquotedblleft{}on the contrary\textquotedblright{}.

\begin{quote}
\ml{അവൾ} \ml{ദേഷ്യപ്പെട്ടില്ല}, \ml{മറിച്ച്} \ml{ചിരിച്ചു}.
\end{quote}

\begin{quote}
\emph{avaḷ dēṣyappeṭṭilla, maṟiccŭ ciriccu.}
\end{quote}

\begin{quote}
She did not get angry; on the contrary, she laughed.
\end{quote}

It follows a negative clause and turns it round.

\begin{cousinweb}
\textbf{\ml{മറിച്ച്}} (\emph{maṟiccŭ}) is native Dravidian, the -\ml{ച്ച്} form of \ml{മറിക്കുക}, \textquotedblleft{}to turn over\textquotedblright{} (DEDR 4761): \textquotedblleft{}turning it over\textquotedblright{}.
\end{cousinweb}

\begin{grammarlens}[title={What you've built}]
One more word for giving reasons and opinions.
\end{grammarlens}

\subsection*{Your turn}
\begin{itemize}
\raggedright
  \item \emph{Say it:} \emph{maṟiccŭ}
  \item \emph{Say it:} \emph{maṟiccŭ}, once more
  \item \emph{Say it:} \emph{avaḷ dēṣyappeṭṭilla, maṟiccŭ ciriccu.} — She did not get angry; on the contrary, she laughed.
  \item \emph{Recall:} say the Malayalam for in reality, then the Malayalam for in short, then say \emph{maṟiccŭ} again
\end{itemize}

\begin{tcolorbox}[breakable,colback=teal!4,colframe=teal!35!black,title={Before you move on}]
What does \ml{മറിച്ച്} mean? (\textquotedblleft{}On the contrary\textquotedblright{}.) Say it once more.
\end{tcolorbox}
\section[\texorpdfstring{\ml{അഥവാ} (athavā)}{അഥവാ (athavā)}]{\ml{അഥവാ} (athavā) — or; in other words}

\label{lesson:ML-C412-athava}



\begin{quote}
Before the new one: say the Malayalam for in short, then the Malayalam for on the contrary.
\end{quote}

\subsection*{You'll want to know: \ml{അഥവാ}}
\textbf{\ml{അഥവാ}} — \emph{athavā} — \textquotedblleft{}or; in other words\textquotedblright{}.

\begin{quote}
\ml{ചായ} \ml{അഥവാ} \ml{കാപ്പി}, \ml{ഏതെങ്കിലും} \ml{മതി}.
\end{quote}

\begin{quote}
\emph{cāya athavā kāppi, ēteṅkiluṁ mati.}
\end{quote}

\begin{quote}
Tea or coffee, either will do.
\end{quote}

A bookish or. It can also mean if by chance: \ml{അഥവാ} \ml{അവൻ} \ml{വന്നാൽ}, if he happens to come.

\begin{cousinweb}
\textbf{\ml{അഥവാ}} (\emph{athavā}) is borrowed from Sanskrit \emph{athavā}, \textquotedblleft{}or\textquotedblright{} (CDIAL 231), from \emph{atha}, \textquotedblleft{}then\textquotedblright{}, and \emph{vā}, \textquotedblleft{}or\textquotedblright{}.
\end{cousinweb}

\begin{grammarlens}[title={What you've built}]
One more word for giving reasons and opinions.
\end{grammarlens}

\subsection*{Your turn}
\begin{itemize}
\raggedright
  \item \emph{Say it:} \emph{athavā}
  \item \emph{Say it:} \emph{athavā}, once more
  \item \emph{Say it:} \emph{cāya athavā kāppi, ēteṅkiluṁ mati.} — Tea or coffee, either will do.
  \item \emph{Recall:} say the Malayalam for in short, then the Malayalam for on the contrary, then say \emph{athavā} again
\end{itemize}

\begin{tcolorbox}[breakable,colback=teal!4,colframe=teal!35!black,title={Before you move on}]
What does \ml{അഥവാ} mean? (\textquotedblleft{}Or; in other words\textquotedblright{}.) Say it once more.
\end{tcolorbox}
\section[\texorpdfstring{\ml{പൊതുവേ} (potuvē)}{പൊതുവേ (potuvē)}]{\ml{പൊതുവേ} (potuvē) — generally}

\label{lesson:ML-C412-potuve}



\begin{quote}
Before the new one: say the Malayalam for on the contrary, then the Malayalam for or.
\end{quote}

\subsection*{You'll want to know: \ml{പൊതുവേ}}
\textbf{\ml{പൊതുവേ}} — \emph{potuvē} — \textquotedblleft{}generally\textquotedblright{}.

\begin{quote}
\ml{പൊതുവേ} \ml{കുട്ടികൾക്ക്} \ml{മധുരം} \ml{ഇഷ്ടമാണ്}.
\end{quote}

\begin{quote}
\emph{potuvē kuṭṭikaḷkkŭ madhuraṁ iṣṭamāṇŭ.}
\end{quote}

\begin{quote}
Children generally like sweets.
\end{quote}

The ending -\ml{ഏ} makes an adverb: \ml{പൊതു}, common, \ml{പൊതുവേ}, in general.

\begin{cousinweb}
\textbf{\ml{പൊതുവേ}} (\emph{potuvē}) is native Dravidian: \ml{പൊതു}, \textquotedblleft{}common, general\textquotedblright{} (DEDR 4507), with the adverb ending -\ml{ഏ}.
\end{cousinweb}

\begin{grammarlens}[title={What you've built}]
One more word for giving reasons and opinions.
\end{grammarlens}

\subsection*{Your turn}
\begin{itemize}
\raggedright
  \item \emph{Say it:} \emph{potuvē}
  \item \emph{Say it:} \emph{potuvē}, once more
  \item \emph{Say it:} \emph{potuvē kuṭṭikaḷkkŭ madhuraṁ iṣṭamāṇŭ.} — Children generally like sweets.
  \item \emph{Recall:} say the Malayalam for on the contrary, then the Malayalam for or, then say \emph{potuvē} again
\end{itemize}

\begin{tcolorbox}[breakable,colback=teal!4,colframe=teal!35!black,title={Before you move on}]
What does \ml{പൊതുവേ} mean? (\textquotedblleft{}Generally\textquotedblright{}.) Say it once more.
\end{tcolorbox}
\section[\texorpdfstring{\ml{എന്നുവെച്ചാൽ} (ennuveccāl)}{എന്നുവെച്ചാൽ (ennuveccāl)}]{\ml{എന്നുവെച്ചാൽ} (ennuveccāl) — that means; that is}

\label{lesson:ML-C412-ennuveccal}



\begin{quote}
Before the new one: say the Malayalam for or, then the Malayalam for generally.
\end{quote}

\subsection*{You'll want to know: \ml{എന്നുവെച്ചാൽ}}
\textbf{\ml{എന്നുവെച്ചാൽ}} — \emph{ennuveccāl} — \textquotedblleft{}that means; that is\textquotedblright{}.

\begin{quote}
\ml{നാളെ} \ml{അവധിയാണ്}, \ml{എന്നുവെച്ചാൽ} \ml{ക്ലാസ്} \ml{ഇല്ല}.
\end{quote}

\begin{quote}
\emph{nāḷe avadhiyāṇŭ, ennuveccāl klāsŭ illa.}
\end{quote}

\begin{quote}
Tomorrow is a holiday, which means there is no class.
\end{quote}

In speech often \ml{എന്നു} \ml{പറഞ്ഞാൽ}, if one says so.

\begin{cousinweb}
\textbf{\ml{എന്നുവെച്ചാൽ}} (\emph{ennuveccāl}) is a native Malayalam compound: \ml{എന്ന്}, \textquotedblleft{}thus (said)\textquotedblright{} (DEDR 868), and \ml{വെച്ചാൽ}, \textquotedblleft{}if one puts\textquotedblright{}, from \ml{വെക്കുക}, \textquotedblleft{}to put\textquotedblright{} (DEDR 5549): \textquotedblleft{}if one puts it that way\textquotedblright{}.
\end{cousinweb}

\begin{grammarlens}[title={What you've built}]
One more word for giving reasons and opinions.
\end{grammarlens}

\subsection*{Your turn}
\begin{itemize}
\raggedright
  \item \emph{Say it:} \emph{ennuveccāl}
  \item \emph{Say it:} \emph{ennuveccāl}, once more
  \item \emph{Say it:} \emph{nāḷe avadhiyāṇŭ, ennuveccāl klāsŭ illa.} — Tomorrow is a holiday, which means there is no class.
  \item \emph{Recall:} say the Malayalam for or, then the Malayalam for generally, then say \emph{ennuveccāl} again
\end{itemize}

\begin{tcolorbox}[breakable,colback=teal!4,colframe=teal!35!black,title={Before you move on}]
What does \ml{എന്നുവെച്ചാൽ} mean? (\textquotedblleft{}That means; that is\textquotedblright{}.) Say it once more.
\end{tcolorbox}
\section[(dialogue)]{Words for giving reasons}

\label{lesson:ML-R412-review-405-412}



\begin{quote}
Say the Malayalam for to plead a case and logic.
\end{quote}

\subsection*{Your turn}
Say these aloud:

\begin{itemize}
\raggedright
  \item to plead a case, logic, a fact, to admit, a conclusion
  \item to disagree, a dispute, a debate, to convince, an outlook
  \item to prove, a query, to assume, reality, to imply
  \item to take into account, to evaluate, a stance, an attitude, to justify
  \item to deny, to claim, to allege, impartial, an explanation
  \item to indicate, to emphasise, to highlight, relevant, irrelevant
  \item therefore, hence, because of, actually, in reality
  \item in short, on the contrary, or, generally, that means
\end{itemize}

\begin{tcolorbox}[breakable,colback=teal!4,colframe=teal!35!black,title={Before you move on}]
To plead a case? (\textbf{\ml{വാദിക്കുക}}, \emph{vādikkuka}.) Logic? (\textbf{\ml{യുക്തി}}, \emph{yukti}.) A fact? (\textbf{\ml{വസ്തുത}}, \emph{vastuta}.) To admit? (\textbf{\ml{സമ്മതിക്കുക}}, \emph{sammatikkuka}.) A conclusion? (\textbf{\ml{നിഗമനം}}, \emph{nigamanaṁ}.) To disagree? (\textbf{\ml{വിയോജിക്കുക}}, \emph{viyōjikkuka}.) A dispute? (\textbf{\ml{തർക്കം}}, \emph{tarkkaṁ}.) A debate? (\textbf{\ml{സംവാദം}}, \emph{saṁvādaṁ}.) To convince? (\textbf{\ml{ബോധ്യപ്പെടുത്തുക}}, \emph{bōdhyappeṭuttuka}.) An outlook? (\textbf{\ml{വീക്ഷണം}}, \emph{vīkṣaṇaṁ}.) To prove? (\textbf{\ml{തെളിയിക്കുക}}, \emph{teḷiyikkuka}.) A query? (\textbf{\ml{സംശയം}}, \emph{saṁśayaṁ}.) To assume? (\textbf{\ml{അനുമാനിക്കുക}}, \emph{anumānikkuka}.) Reality? (\textbf{\ml{യാഥാർത്ഥ്യം}}, \emph{yāthārtthyaṁ}.) To imply? (\textbf{\ml{അർത്ഥമാക്കുക}}, \emph{artthamākkuka}.) To take into account? (\textbf{\ml{പരിഗണിക്കുക}}, \emph{parigaṇikkuka}.) To evaluate? (\textbf{\ml{വിലയിരുത്തുക}}, \emph{vilayiruttuka}.) A stance? (\textbf{\ml{നിലപാട്}}, \emph{nilapāṭŭ}.) An attitude? (\textbf{\ml{മനോഭാവം}}, \emph{manōbhāvaṁ}.) To justify? (\textbf{\ml{ന്യായീകരിക്കുക}}, \emph{nyāyīkarikkuka}.) To deny? (\textbf{\ml{നിഷേധിക്കുക}}, \emph{niṣēdhikkuka}.) To claim? (\textbf{\ml{അവകാശപ്പെടുക}}, \emph{avakāśappeṭuka}.) To allege? (\textbf{\ml{ആരോപിക്കുക}}, \emph{ārōpikkuka}.) Impartial? (\textbf{\ml{നിഷ്പക്ഷമായ}}, \emph{niṣpakṣamāya}.) An explanation? (\textbf{\ml{വിശദീകരണം}}, \emph{viśadīkaraṇaṁ}.) To indicate? (\textbf{\ml{സൂചിപ്പിക്കുക}}, \emph{sūcippikkuka}.) To emphasise? (\textbf{\ml{ഊന്നിപ്പറയുക}}, \emph{ūnnippaṟayuka}.) To highlight? (\textbf{\ml{ചൂണ്ടിക്കാണിക്കുക}}, \emph{cūṇṭikkāṇikkuka}.) Relevant? (\textbf{\ml{പ്രസക്തമായ}}, \emph{prasaktamāya}.) Irrelevant? (\textbf{\ml{അപ്രസക്തമായ}}, \emph{aprasaktamāya}.) Therefore? (\textbf{\ml{അതിനാൽ}}, \emph{atināl}.) Hence? (\textbf{\ml{ആയതിനാൽ}}, \emph{āyatināl}.) Because of? (\textbf{\ml{മൂലം}}, \emph{mūlaṁ}.) Actually? (\textbf{\ml{യഥാർത്ഥത്തിൽ}}, \emph{yathārtthattil}.) In reality? (\textbf{\ml{വാസ്തവത്തിൽ}}, \emph{vāstavattil}.) In short? (\textbf{\ml{ചുരുക്കത്തിൽ}}, \emph{curukkattil}.) On the contrary? (\textbf{\ml{മറിച്ച്}}, \emph{maṟiccŭ}.) Or? (\textbf{\ml{അഥവാ}}, \emph{athavā}.) Generally? (\textbf{\ml{പൊതുവേ}}, \emph{potuvē}.) That means? (\textbf{\ml{എന്നുവെച്ചാൽ}}, \emph{ennuveccāl}.)
\end{tcolorbox}
